% theorem1_revised.tex
% Revised Theorem 1: Well-posedness of the ELAPSE SDE
% Non-circular proof: local existence -> non-negativity -> global existence
%                     -> uniqueness (in strict logical order).
%
% The original proof had a circular dependency: the Lyapunov argument
% (global existence) invoked Markov's inequality on M(t)>0, which was
% not yet established since non-negativity was proved later using global
% existence. This revision removes that circularity via the Feller
% boundary classification.
%
% References:
%   [Ok03] Oksendal (2003) SDE 6th ed., Theorem 5.2.1 (Picard-Lindelof)
%   [RY99] Revuz-Yor (1999) Brownian Motion, Theorem IX.3.7 (comparison)
%   [Fe51] Feller (1951) Ann. Math. 54, pp. 173-182 (boundary classification)
%   [YW71] Yamada-Watanabe (1971) (pathwise uniqueness for CIR)

\begin{theorem}[Well-posedness of the ELAPSE SDE]
\label{thm:wellposed}
Fix $T>0$, $\alpha,\mu,\sigma_n>0$, source $s\ge 0$, and ensemble weight
$w\in\mathcal{W}$.  Suppose $x_0\in\mathbb{R}^n_+$ (componentwise non-negative)
and $x_{0,i}>0$ for at least one index $i$.  Then the ELAPSE SDE
\begin{equation}
  \mathrm{d}x_i = \underbrace{\bigl[-\alpha(Lx)_i - \mu\,\Delta(t;w)\,x_i + s_i\bigr]}_{b_i(x,t)}\mathrm{d}t
                  + \sigma_n\sqrt{x_i}\,\mathrm{d}W_i, \quad i=1,\ldots,n,
  \label{eq:elapse_sde_thm1}
\end{equation}
with independent Brownian motions $(W_i)$ on a filtered probability space,
admits a unique strong solution $(x(t))_{t\in[0,T]}$ satisfying:
\begin{enumerate}[label=(\roman*)]
  \item \textbf{Local existence:} there exists a stopping time $\tau_K>0$ a.s.\
    such that $x(t)$ is well-defined on $[0,\tau_K)$ for each compact
    $K\subset\mathbb{R}^n_+$.
  \item \textbf{Non-negativity:} $x_i(t)\ge 0$ for all $i$ and all
    $t\in[0,\tau_K)$ a.s., without requiring global existence.
  \item \textbf{Global existence:} the maximal solution interval is $[0,T]$;
    i.e.\ $P(\tau_K<T)\to 0$ as $K\nearrow\mathbb{R}^n_+$.
  \item \textbf{Pathwise uniqueness:} if $(x(t))$ and $(\tilde x(t))$ are
    two solutions with $x(0)=\tilde x(0)=x_0$, then
    $P(x(t)=\tilde x(t)\;\forall t\in[0,T])=1$.
\end{enumerate}
\end{theorem}

\begin{proof}
We establish the four parts \emph{in strict logical order}, avoiding circular
dependencies.

\medskip
\noindent\textbf{Step A: Local existence (part i).}
On any compact $K=[-R,R]^n\cap\mathbb{R}^n_+$, the drift $b_i(x,t)$ is bounded
and globally Lipschitz in $x$ (the Laplacian term is linear; the mortality term
$-\mu\Delta(t)x_i$ is Lipschitz in $x_i$ since $\Delta\in[0,1]$; the source
$s_i\ge 0$ is constant).
The diffusion coefficient $\sigma_i(x)=\sigma_n\sqrt{x_i}$ satisfies
$|\sigma_i(x)-\sigma_i(y)|\le\sigma_n|x_i-y_i|^{1/2}$, which is
$1/2$-H\"older continuous, sufficient for strong existence by the
Yamada--Watanabe criterion \cite{YamadaWatanabe1971}.
By the SDE Picard--Lindel\"of theorem \cite[Thm.~5.2.1]{Oksendal2003},
there exists a unique local strong solution on $[0,\tau_K)$ where
$\tau_K=\inf\{t:\|x(t)\|>R\}$ is the first exit time from $K$.

\medskip
\noindent\textbf{Step B: Non-negativity on $[0,\tau_K)$ (part ii).}
We establish non-negativity on the local solution interval \emph{without}
appealing to global existence.
Consider a single component $x_i$ on $[0,\tau_K)$.
At the boundary $x_i=0$, the drift satisfies
\begin{equation}
  b_i(x,t)\big|_{x_i=0}
  = \alpha\sum_{j\sim i} A_{ij}\,x_j + s_i
  \;\ge\; 0,
  \label{eq:inward_drift}
\end{equation}
since $A_{ij}\ge 0$ (adjacency matrix entries), $x_j\ge 0$ on $[0,\tau_K)$
(by induction on the iteration/localisation), and $s_i\ge 0$.
The diffusion coefficient vanishes: $\sigma_n\sqrt{0}=0$.
By Feller's boundary classification \cite{Feller1951}, the combination of
non-negative inward drift and vanishing diffusion at $x_i=0$ makes the
boundary \emph{unattainable} from above: $x_i(t)$ cannot reach $0$ in finite
time if $x_i(0)\ge 0$.  Formally, apply the comparison theorem
\cite[Thm.~IX.3.7]{RevuzYor1999}: let $Y_i$ satisfy
$\mathrm{d}Y_i=\sigma_n\sqrt{Y_i}\,\mathrm{d}W_i$ (zero-drift CIR boundary process)
with $Y_i(0)=0$.  Since $b_i|_{x_i=0}\ge 0=$ (drift of $Y_i$), the comparison
theorem gives $x_i(t)\ge Y_i(t)=0$ a.s.\ on $[0,\tau_K)$.

Hence $x_i(t)\ge 0$ for all $i$ on $[0,\tau_K)$, established \emph{independently}
of global existence.

\medskip
\noindent\textbf{Step C: Global existence (part iii).}
\emph{Now} that non-negativity is established on $[0,\tau_K)$, we have
$M(t)=\sum_i x_i(t)\ge 0$ on $[0,\tau_K)$.
Apply Ito's formula to the Lyapunov function $V(x)=M(x)=\mathbf{1}^\top x$:
\begin{equation}
  \mathrm{d}M = \bigl[-\mu\,\Delta(t)\,M + \|s\|_1\bigr]\,\mathrm{d}t
                + \sigma_n\sum_i\sqrt{x_i}\,\mathrm{d}W_i,
  \label{eq:dM_lyapunov}
\end{equation}
using $\mathbf{1}^\top L=0$ (row-sum zero Laplacian).
Taking expectations and using $\Delta(t)\ge 0$:
\begin{equation}
  \frac{\mathrm{d}}{\mathrm{d}t}\,\mathbb{E}[M(t)] \;\le\; \|s\|_1,
  \quad\Rightarrow\quad
  \mathbb{E}[M(t)] \;\le\; M_0 + \|s\|_1\,t \;<\; \infty \quad \forall\,t\in[0,T].
  \label{eq:EM_bound}
\end{equation}
By Markov's inequality, $P(M(t)>R)\le(M_0+\|s\|_1 T)/R\to 0$ as $R\to\infty$.
Combined with Doob's maximal inequality applied to the martingale part
and a standard localisation argument (see \cite[Thm.~5.2.1]{Oksendal2003}),
$P(\tau_K<T)\to 0$ as $K\nearrow\mathbb{R}^n_+$.
Hence the solution exists globally on $[0,T]$.

\medskip
\noindent\textbf{Step D: Pathwise uniqueness (part iv).}
Let $x(t)$ and $\tilde{x}(t)$ be two solutions with the same initial condition.
Set $e_i=x_i-\tilde{x}_i$.  By the Yamada--Watanabe pathwise uniqueness
criterion \cite{YamadaWatanabe1971}, it suffices to show that the diffusion
coefficient $\sigma_i(x)=\sigma_n\sqrt{x_i}$ satisfies
$|\sigma_i(x)-\sigma_i(y)|\le\rho(|x_i-y_i|)$ for a concave increasing
$\rho$ with $\int_0^1 \rho^{-2}(u)\,\mathrm{d}u=\infty$.
Taking $\rho(u)=\sigma_n\sqrt{u}$ gives $\int_0^1 u^{-1}\,\mathrm{d}u=+\infty$,
satisfying the Yamada--Watanabe criterion.
Combined with the Lipschitz drift established in Step~A, pathwise uniqueness follows.
\end{proof}

\begin{remark}[Logical ordering]
\label{rem:thm1_ordering}
The proof order---Steps A, B, C, D---is essential.  Non-negativity (Step~B) is
established using only the local solution from Step~A, specifically via the
Feller boundary argument that requires neither $M(t)>0$ nor global existence.
The Lyapunov/Gronwall argument (Step~C) then \emph{uses} $M(t)\ge 0$ (proved in B)
to conclude that $\mathbb{E}[M(t)]<\infty$, and hence that the solution cannot
blow up.  Any proof that invokes $M(t)>0$ \emph{before} establishing
non-negativity would be circular.
\end{remark}
